\chapter{What We're Forecasting}
\label{ch:what-we-forecast}

Realized variance on day $t$ is $\RV_t = \sum_{i=1}^{M} r_{t,i}^2$, the sum of $M$ squared intraday returns sampled at 5-minute frequency \citep{ABDL2003}.
Our target is $\log \RV_{t+h}$ for horizons $h \in \{1, 5, 22\}$ trading days, where multi-day targets aggregate: $\RV_{t+1:t+h} = h^{-1}\sum_{j=1}^{h} \RV_{t+j}$.

%% ──────────────────────────────────────────────
\section{The Universe}
%% ──────────────────────────────────────────────

The model is trained and evaluated on 35 instruments over 11.3 years of history (Jan 2012 to Mar 2023).

\begin{table}[ht]
\centering
\caption{Forecast universe: 35 instruments.}
\label{tab:universe}
\begin{tabular}{@{}lll@{}}
\toprule
Category & Count & Examples \\
\midrule
Mega-cap US equities & 30 & AAPL, MSFT, JPM, JNJ, XOM, AMZN \\
Broad-market ETFs    & 4  & SPY, QQQ, IWM, DIA \\
Equity index futures & 1  & E-mini S\&P 500 (ES) \\
\midrule
\textbf{Total}       & \textbf{35} & \\
\bottomrule
\end{tabular}
\end{table}

All RV series are computed from 5-minute returns using the realized kernel estimator of \citet{BNHLS2008}, which corrects for microstructure noise.
Models are fit per-asset (no pooling across names).

%% ──────────────────────────────────────────────
\section{Success Criteria}
%% ──────────────────────────────────────────────

The primary metric is QLIKE, a quasi-likelihood loss that penalizes relative forecast errors and is robust to noisy volatility proxies \citep{Patton2011}.
The formal definition appears in Chapter~13.

\begin{table}[ht]
\centering
\caption{Evaluation criteria.}
\label{tab:criteria}
\begin{tabular}{@{}llp{7.5cm}@{}}
\toprule
Priority & Metric & Requirement \\
\midrule
Primary   & QLIKE & 30--80 bps improvement over HARQ baseline \citep{BPQ2016}, averaged across universe \\
Secondary & MSE, MAE & Reported alongside QLIKE for robustness \\
Secondary & Diebold-Mariano test & Statistically significant improvement ($p < 0.05$) vs.\ each baseline \citep{DieboldMariano1995} \\
Tertiary  & Economic value & Out-of-sample utility gain in a volatility-targeting portfolio \citep{MoreiraMuir2017} \\
\bottomrule
\end{tabular}
\end{table}

%% ──────────────────────────────────────────────
\section{The High-Level Pipeline}
%% ──────────────────────────────────────────────

Figure~\ref{fig:pipeline} shows the end-to-end forecasting system.
Raw data from six sources feeds into a layered feature-engineering pipeline (Layers 0--7, detailed in Chapters 3--8).
Two model families consume the resulting feature matrix: LightGBM (Ch.~9) and an LSTM operating on intraday return sequences (Ch.~10).
Their forecasts are combined via an ensemble blend (Ch.~11) and evaluated against QLIKE.

\begin{figure}[ht]
\centering
\begin{tikzpicture}[node distance=1.8cm, every node/.style={font=\small}]
  % Nodes
  \node[flowblock, minimum width=2.8cm] (data) {Raw Data\\(6 sources)};
  \node[flowblock, minimum width=2.8cm, right=of data] (feat) {Feature Engineering\\(Layers 0--7)};
  \node[flowblock, minimum width=2.2cm, above right=0.6cm and 1.8cm of feat] (lgbm) {LightGBM};
  \node[flowblock, minimum width=2.2cm, below right=0.6cm and 1.8cm of feat] (lstm) {LSTM};
  \node[flowblock, minimum width=2.8cm, right=2.8cm of feat] (ens) {Ensemble\\Blend};
  \node[flowblock, minimum width=2.8cm, right=of ens] (out) {Forecast:\\$\log \RV_{t+h}$};
  \node[flowblock, minimum width=2.2cm, right=of out] (eval) {Evaluation\\(QLIKE)};

  % Arrows
  \draw[arrow] (data) -- (feat);
  \draw[arrow] (feat) -- (lgbm);
  \draw[arrow] (feat) -- (lstm);
  \draw[arrow] (lgbm) -- (ens);
  \draw[arrow] (lstm) -- (ens);
  \draw[arrow] (ens) -- (out);
  \draw[arrow] (out) -- (eval);
\end{tikzpicture}
\caption{High-level forecasting pipeline.}
\label{fig:pipeline}
\end{figure}

The six raw data sources are: (i) tick-level trade prices, (ii) daily OHLCV, (iii) E-mini S\&P 500 Level~2 order book, (iv) SPX implied volatility surface, (v) cross-asset prices (VIX, bonds, gold, oil, FX), and (vi) earnings and macro event calendars.

\begin{keyidea}[Features Over Models]
The feature set matters more than the model.
Layers 0--2 (approximately 20 features covering HAR components, jump decomposition, and signed volatility) achieve 85\% of the forecasting accuracy attainable with the full 120-feature set.
\end{keyidea}
